%=============================================================================
\section{Pipeline Overview}
\label{sec:pipeline}
%=============================================================================

A multi-AIR computation produces $A$ independent STARK proofs,
one per AIR instance.
The aggregation pipeline transforms these into a single proof
through up to seven stages:

\begin{center}
\renewcommand{\arraystretch}{1.3}
\begin{tabular}{@{}clll@{}}
  \toprule
  Stage & Name & Input & Output \\
  \midrule
  1 & Basic          & AIR witness           & Per-AIR STARK proof \\
  2 & Compressor     & Basic STARK proof     & Compressed STARK proof \\
  3 & Recursive1     & Basic/Compressor      & Normalized SNARK proof \\
  4 & Recursive2     & $\leq 3$ Recursive1/2 & Aggregated SNARK proof \\
  5 & VadcopFinal    & One R2 per airgroup   & Cross-AIR verified STARK proof \\
  \midrule
  \multicolumn{4}{@{}l}{\itshape The following stages are not currently deployed
    (see note below):} \\
  6 & RecursiveF     & VadcopFinal           & BN254-native STARK proof \\
  7 & fflonk         & RecursiveF            & Ethereum-verifiable SNARK \\
  \bottomrule
\end{tabular}
\end{center}

\medskip

Stages 1--5 operate entirely over the Goldilocks field $\F$ with Poseidon2 hashing.
Stage 6 would transition to BN254 with Poseidon hashing over the BN254 scalar field.
Stage 7 would produce a pairing-based SNARK verifiable by an on-chain Solidity contract.

\paragraph{Current deployment.}
In the current Zisk ethproofs integration, the pipeline terminates at
Stage~5 (VadcopFinal).
The coordinator sets \texttt{final\_snark\,=\,false}, which skips
Stages~6 and~7 entirely.
The VadcopFinal STARK proof---a Goldilocks-native proof with no
elliptic curve cryptography---is submitted directly to the
ethproofs aggregation service.
Stages~6 and~7 (RecursiveF and fflonk) are implemented but reserved
for a future configuration where on-chain STARK verification
is replaced by a constant-size pairing-based SNARK.

\bigskip

\begin{figure}[h]
\centering
\begin{tikzpicture}[
    % Node styles
    basic/.style={draw, rounded corners=2pt, minimum width=1.4cm, minimum height=0.7cm,
                  font=\footnotesize, fill=blue!8},
    comp/.style={draw, rounded corners=2pt, minimum width=1.4cm, minimum height=0.7cm,
                 font=\footnotesize, fill=orange!12},
    r1/.style={draw, rounded corners=2pt, minimum width=1.4cm, minimum height=0.7cm,
               font=\footnotesize, fill=green!10},
    r2/.style={draw, rounded corners=2pt, minimum width=1.8cm, minimum height=0.7cm,
               font=\footnotesize, fill=violet!10},
    final/.style={draw, rounded corners=3pt, minimum width=2.2cm, minimum height=0.8cm,
                  font=\footnotesize\bfseries, fill=red!8},
    stage/.style={font=\footnotesize\itshape, text=black!60},
    arr/.style={-{Stealth[length=4pt]}, thick, black!70},
    dotarr/.style={-{Stealth[length=4pt]}, thick, black!70, densely dashed},
  ]

  % --- Stage labels (left margin) ---
  \node[stage] at (-4.8, 0)    {Stage 1};
  \node[stage] at (-4.8,-1.5)  {Stage 2};
  \node[stage] at (-4.8,-3.0)  {Stage 3};
  \node[stage] at (-4.8,-4.8)  {Stage 4};
  \node[stage] at (-4.8,-7.2)  {Stage 5};
  \node[stage] at (-4.8,-8.7)  {Stage 6};
  \node[stage] at (-4.8,-10.2) {Stage 7};

  % --- Stage 1: Basic proofs ---
  % Small AIRs (no compressor needed)
  \node[basic] (b1) at (-3.2, 0) {Main};
  \node[basic] (b2) at (-1.6, 0) {Rom};
  \node[basic] (b3) at ( 0.0, 0) {Mem};
  \node[basic] (b4) at ( 1.6, 0) {Binary};
  % Large AIRs (need compressor)
  \node[basic] (b5) at ( 3.2, 0) {Keccakf};
  \node[basic, minimum width=0.8cm] (bdots) at (4.5, 0) {$\cdots$};

  % --- Stage 2: Compressor (only for large AIRs) ---
  \node[comp] (c5) at (3.2, -1.5) {Compr.};

  % Arrows: small AIRs skip compressor
  \draw[arr] (b1) -- ++(0,-0.6) -| ([xshift=-0.4cm]b1.south |- 0,-3.0);
  \draw[arr] (b2) -- ++(0,-0.6) -| ([xshift=-0.4cm]b2.south |- 0,-3.0);
  \draw[arr] (b3) -- ++(0,-0.6) -| ([xshift=-0.4cm]b3.south |- 0,-3.0);
  \draw[arr] (b4) -- ++(0,-0.6) -| ([xshift=-0.4cm]b4.south |- 0,-3.0);
  % Arrow: large AIR goes through compressor
  \draw[arr] (b5) -- (c5);
  \draw[arr] (c5) -- ++(0,-0.6) -| ([xshift=-0.4cm]c5.south |- 0,-3.0);

  % "skip" annotation
  \node[font=\tiny, text=black!50] at (-0.8, -1.5) {(skip if verifier $\leq 2^{17}$ rows)};

  % --- Stage 3: Recursive1 ---
  \node[r1] (r1a) at (-3.2, -3.0) {R1};
  \node[r1] (r1b) at (-1.6, -3.0) {R1};
  \node[r1] (r1c) at ( 0.0, -3.0) {R1};
  \node[r1] (r1d) at ( 1.6, -3.0) {R1};
  \node[r1] (r1e) at ( 3.2, -3.0) {R1};
  \node[r1, minimum width=0.8cm] (r1dots) at (4.5, -3.0) {$\cdots$};

  % --- Stage 4: Recursive2 tree ---
  % First level: groups of 3
  \node[r2] (r2a) at (-1.6, -4.5) {Recursive2};
  \node[r2] (r2b) at ( 1.6, -4.5) {Recursive2};

  \draw[arr] (r1a) -- (r2a);
  \draw[arr] (r1b) -- (r2a);
  \draw[arr] (r1c) -- (r2a);
  \draw[arr] (r1d) -- (r2b);
  \draw[arr] (r1e) -- (r2b);
  \draw[dotarr] (r1dots) -- (r2b);

  % Second level: merge
  \node[r2] (r2final) at (0, -6.3) {Recursive2};

  \draw[arr] (r2a) -- (r2final);
  \draw[arr] (r2b) -- (r2final);

  % "3-to-1 tree" annotation
  \node[font=\tiny, text=black!50, anchor=west] at (3.0, -5.4) {$\lceil\log_3 A\rceil$ levels};

  % Airgroup brace
  \begin{scope}[on background layer]
    \node[draw=black!20, fill=black!3, rounded corners=4pt,
          fit=(b1)(bdots)(r2final), inner sep=6pt] (airgroup) {};
  \end{scope}
  \node[font=\scriptsize, text=black!50, anchor=south] at (airgroup.north) {Airgroup ``Zisk'' (21 AIRs)};

  % --- Stage 5: VadcopFinal ---
  \node[final] (vf) at (0, -7.8) {VadcopFinal};
  \draw[arr] (r2final) -- (vf);
  \node[font=\tiny, text=black!50, anchor=west] at (1.5, -7.8) {global challenge + global constraints};

  % --- Stages 6--7: not currently deployed ---
  \draw[densely dashed, black!30] (-4.0, -8.6) -- (5.0, -8.6);
  \node[font=\tiny\itshape, text=black!40, anchor=west] at (-4.0, -8.85)
    {not currently deployed};

  % --- Stage 6: RecursiveF ---
  \node[final, fill=yellow!10, draw=black!30, text=black!40] (rf) at (0, -9.3) {RecursiveF};
  \draw[dotarr, black!30] (vf) -- (rf);
  \node[font=\tiny, text=black!35, anchor=west] at (1.5, -9.3) {Goldilocks $\to$ BN254};

  % --- Stage 7: fflonk ---
  \node[final, fill=green!8, draw=black!30, text=black!40] (ff) at (0, -10.8) {fflonk};
  \draw[dotarr, black!30] (rf) -- (ff);
  \node[font=\tiny, text=black!35, anchor=west] at (1.5, -10.8) {Ethereum-verifiable};

\end{tikzpicture}
\caption{VADCOP recursive aggregation pipeline.
  Each per-AIR Basic proof is optionally compressed,
  normalized by Recursive1,
  then aggregated via a Recursive2 tree within each airgroup.
  VadcopFinal verifies the global challenge and cross-AIR constraints.
  Stages below the dashed line (RecursiveF, fflonk) are not currently deployed.}
\label{fig:pipeline}
\end{figure}

%-----------------------------------------------------------------------------
\subsection{Proof Type Table}
\label{sec:proof-types}

\begin{center}
\renewcommand{\arraystretch}{1.3}
\begin{tabular}{@{}llp{0.48\textwidth}@{}}
  \toprule
  Type & System & Description \\
  \midrule
  Basic & STARK & Per-AIR STARK proof
    (see the \emph{STARK Protocol}) \\
  Compressor & STARK & Optional compression for large Basic proofs \\
  Recursive1 & Circom SNARK & Wraps one Basic/Compressor STARK
                                in a Groth16-style circuit \\
  Recursive2 & Circom SNARK & Aggregates $\leq 3$ proofs of the same airgroup
                                into 1 (tree reduction) \\
  VadcopFinal & STARK & Combines one Recursive2 per airgroup;
                          verifies global challenge and global constraints \\
  RecursiveF & STARK & Re-proves VadcopFinal over BN254 for pairing-friendly verification \\
  fflonk & SNARK & Ethereum-verifiable proof (fflonk variant of PLONK) \\
  \bottomrule
\end{tabular}
\end{center}

%-----------------------------------------------------------------------------
\subsection{Field Transitions}
\label{sec:field-transitions}

\begin{itemize}[nosep]
  \item \textbf{Stages 1--5} (Basic through VadcopFinal):
        Goldilocks field $\F = \mathbb{F}_{2^{64} - 2^{32} + 1}$,
        Poseidon2 hashing, cubic extension $\Fext$.
        In the current deployment, the entire pipeline uses this single field.
  \item \textbf{Stage 6} (RecursiveF, not deployed):
        BN254 scalar field $\mathbb{F}_r$ where
        $r$ is the order of the BN254 curve group.
        Poseidon hashing over $\mathbb{F}_r$.
        Goldilocks arithmetic is emulated within BN254 circuits.
  \item \textbf{Stage 7} (fflonk, not deployed):
        BN254 pairing-based SNARK.
        The output would be a constant-size proof verifiable by
        a Solidity smart contract.
\end{itemize}

%=============================================================================
\section{Basic STARK Proof (Stage 1)}
\label{sec:basic}
%=============================================================================

Each AIR instance produces an independent STARK proof using the protocol
described in the \emph{STARK Protocol}.

\begin{enumerate}[label=(\roman*)]
  \item The prover commits stage-1 polynomials and computes a
        contribution $\kappa^{(a)}$ for challenge aggregation
        (see the \emph{STARK Protocol}, Multi-AIR Challenge Binding).

  \item After the global challenge $\chi$ is derived from all contributions
        (see the \emph{STARK Protocol}, Challenge Aggregation),
        each AIR seeds its transcript with $\chi$ and runs the full
        commitment and query phases.

  \item The output is a binary proof $\pi^{(a)}_{\mathrm{Basic}}$
        containing Merkle roots, polynomial evaluations,
        FRI layers, grinding nonce, and Merkle opening proofs.
\end{enumerate}

%=============================================================================
\section{Compressor (Stage 2, Optional)}
\label{sec:compressor}
%=============================================================================

The Compressor is an optional stage applied when the STARK verifier circuit
for a particular AIR exceeds $2^{17}$ rows.

\begin{enumerate}[label=(\roman*)]
  \item \textbf{Circuit.}
        A Circom circuit over the Goldilocks field that verifies one Basic STARK proof.
        The circuit's execution trace forms a new AIR.

  \item \textbf{When needed.}
        If the number of verifier circuit rows is $\leq 2^{17}$,
        the Compressor is skipped and the Basic proof proceeds directly
        to Recursive1.

  \item \textbf{Output.}
        A new STARK proof $\pi^{(a)}_{\mathrm{Comp}}$ with a smaller
        verification circuit than the original Basic proof.

  \item \textbf{stage1Hash.}
        The Compressor computes a commitment hash:
        \[
          \mathrm{stage1Hash}^{(a)} = \Poseidon\bigl(\mathrm{vk}^{(a)},\;
          r_1^{(a)},\; \mathrm{airValues}^{(a)}\bigr),
        \]
        which chains through subsequent recursive layers.
\end{enumerate}

%=============================================================================
\section{Recursive1: Per-AIR Normalization (Stage 3)}
\label{sec:recursive1}
%=============================================================================

Recursive1 wraps each Basic or Compressor STARK proof in a SNARK circuit,
producing a normalized proof format suitable for tree aggregation.

\begin{enumerate}[label=(\roman*)]
  \item \textbf{Circuit.}
        A Circom SNARK circuit that verifies one STARK proof.
        The verification key is selected based on the input proof's circuit type:
        \begin{itemize}[nosep]
          \item Circuit type = ``basic'': use the Basic STARK verification key.
          \item Circuit type = ``compressor'': use the Compressor verification key.
        \end{itemize}

  \item \textbf{Input.}
        One STARK proof $\pi^{(a)}$ (either Basic or Compressor).

  \item \textbf{Output.}
        A normalized proof $\pi^{(a)}_{\mathrm{R1}}$ with a standardized
        output format containing:
        \begin{itemize}[nosep]
          \item Public inputs and proof values.
          \item The global challenge $\chi$.
          \item The root of the aggregated challenge computation $\mathrm{rootCAgg}$.
          \item The $\mathrm{stage1Hash}$ for this AIR instance.
        \end{itemize}
\end{enumerate}

%=============================================================================
\section{Recursive2: Tree Aggregation (Stage 4)}
\label{sec:recursive2}
%=============================================================================

Recursive2 aggregates proofs within the same airgroup using a 3-to-1
tree reduction.

%-----------------------------------------------------------------------------
\subsection{Aggregation Structure}
\label{sec:r2-structure}

\begin{enumerate}[label=(\roman*)]
  \item \textbf{Fan-in.}
        Each Recursive2 circuit takes up to
        $N_{\mathrm{agg}} = 3$ input proofs.
        These may be Recursive1 proofs or other Recursive2 proofs
        from the same airgroup.

  \item \textbf{Null-proof padding.}
        If the number of proofs is not divisible by~3,
        null proofs are used as padding.
        The circuit distinguishes null proofs via a circuit type flag
        (circuit type $= 0$ indicates a null proof that is not verified).

  \item \textbf{Verification key selection.}
        The circuit selects the appropriate verification key based on
        the circuit type of each input:
        \begin{itemize}[nosep]
          \item Type 0: null proof (skip verification).
          \item Type 1: aggregated proof (Recursive2 verification key).
          \item Type $\geq 2$: basic proof (Recursive1 verification key,
                with type encoding the specific AIR).
        \end{itemize}
\end{enumerate}

%-----------------------------------------------------------------------------
\subsection{Airgroup Value Aggregation}
\label{sec:r2-airgroup-values}

Each input proof carries airgroup values (accumulated $\mathrm{gsum}$ and
$\mathrm{gprod}$ boundary values).
Recursive2 aggregates these:

\begin{itemize}[nosep]
  \item \textbf{Sum type} ($\mathrm{gsum}$):
        add the values from all non-null inputs.
  \item \textbf{Product type} ($\mathrm{gprod}$):
        multiply the values from all non-null inputs.
\end{itemize}

%-----------------------------------------------------------------------------
\subsection{stage1Hash Chaining}
\label{sec:r2-stage1hash}

The $\mathrm{stage1Hash}$ values from input proofs are chained
using Poseidon2:
\[
  \mathrm{stage1Hash}_{\mathrm{out}}
  = \Poseidon\bigl(\mathrm{stage1Hash}_A \;\|\; \mathrm{stage1Hash}_B\bigr).
\]
For three inputs $A, B, C$:
first combine $A$ and $B$, then combine the result with $C$.
Null proofs contribute a zero hash.

%-----------------------------------------------------------------------------
\subsection{Tree Depth}
\label{sec:r2-depth}

For an airgroup with $N_{\mathrm{AIR}}$ instances, the tree has depth
\[
  D = \lceil \log_3(N_{\mathrm{AIR}}) \rceil.
\]
The output is a single proof $\pi^{(g)}_{\mathrm{R2}}$ per airgroup~$g$.

%=============================================================================
\section{VadcopFinal (Stage 5)}
\label{sec:vadcop-final}
%=============================================================================

The VadcopFinal circuit is itself an AIR whose constraints enforce
cross-airgroup consistency.
It takes one Recursive2 proof per airgroup and produces a single STARK proof.

%-----------------------------------------------------------------------------
\subsection{STARK Verification}
\label{sec:vf-stark}

For each airgroup $g$, the VadcopFinal circuit verifies the Recursive2 proof
$\pi^{(g)}_{\mathrm{R2}}$
using the appropriate verification key
(selected between Basic and aggregation keys
depending on the circuit type flag).

%-----------------------------------------------------------------------------
\subsection{Global Challenge Recomputation}
\label{sec:vf-challenge}

The circuit recomputes the global challenge from the public inputs:
\[
  \chi' = \T_G\bigl(\mathrm{publics},\;
                    \mathrm{proofValues},\;
                    \mathrm{stage1Hash}^{(0)}, \ldots,
                    \mathrm{stage1Hash}^{(G-1)}\bigr),
\]
where each $\mathrm{stage1Hash}^{(g)}$ is the chained Poseidon2 hash
accumulated through the Recursive2 tree for airgroup~$g$.

\textbf{Check:} $\chi' = \chi$ (the global challenge used by all per-AIR proofs).

%-----------------------------------------------------------------------------
\subsection{Global Constraint Verification}
\label{sec:vf-global}

The circuit enforces all cross-airgroup constraints.
For a concrete machine, these are defined by the machine specification
(e.g.\ bus balance equations, continuation anchoring;
see the \emph{Zisk Machine} specification).

In general, global constraints take two forms:

\begin{itemize}[nosep]
  \item \textbf{Sum type (logup):}
        the sum of all $\mathrm{gsum}$ boundary values across AIRs
        sharing a bus equals zero.
        \[
          \sum_{a \in \mathrm{bus}(b)} \mathrm{gsum}^{(a)}[N-1] = 0
          \quad \text{for each bus } b.
        \]

  \item \textbf{Product type (permutation):}
        the product of all $\mathrm{gprod}$ boundary values across AIRs
        sharing a bus equals one.
        \[
          \prod_{a \in \mathrm{bus}(b)} \mathrm{gprod}^{(a)}[N-1] = 1
          \quad \text{for each bus } b.
        \]
\end{itemize}

%-----------------------------------------------------------------------------
\subsection{Output}
\label{sec:vf-output}

The VadcopFinal proof $\pi_{\mathrm{VF}}$ is itself a STARK proof
over the Goldilocks field,
verified by the standard STARK verifier
(see the \emph{STARK Protocol}, Query Phase).
Its transcript is seeded directly with
$(\mathrm{vk},\; \Hash(\mathrm{pub}),\; r_1)$
rather than with a global challenge,
since it \emph{is} the outermost layer.

In the current deployment, this is the \textbf{final proof} submitted to
the ethproofs service.
The proof is serialized as a flat array of Goldilocks field elements
(\texttt{Vec<u64>}) and transmitted as base64-encoded binary.

%=============================================================================
\section{RecursiveF: Field Transition (Stage 6, Not Deployed)}
\label{sec:recursivef}
%=============================================================================

\emph{This stage is implemented but not currently deployed.
The pipeline terminates at VadcopFinal (Stage~5) in the current configuration.}

\medskip

RecursiveF re-proves the VadcopFinal STARK over the BN254 scalar field,
enabling verification by a pairing-based SNARK in the final stage.

\begin{enumerate}[label=(\roman*)]
  \item \textbf{Purpose.}
        The Goldilocks-based STARK proof from VadcopFinal cannot be
        directly verified by a BN254 pairing check.
        RecursiveF ``translates'' the proof to BN254-native arithmetic.

  \item \textbf{BN254-native Merkle hashing.}
        The Merkle trees in the RecursiveF proof use Poseidon over the
        BN254 scalar field $\mathbb{F}_r$ instead of Poseidon2 over Goldilocks.

  \item \textbf{Goldilocks emulation.}
        Goldilocks field arithmetic ($p = 2^{64} - 2^{32} + 1$) is
        emulated within BN254 circuits.
        This is the computational bottleneck of the recursion pipeline.

  \item \textbf{Output.}
        A STARK proof $\pi_{\mathrm{RF}}$ over BN254,
        ready for the fflonk stage.
\end{enumerate}

%=============================================================================
\section{fflonk: Ethereum-Verifiable Proof (Stage 7, Not Deployed)}
\label{sec:fflonk}
%=============================================================================

\emph{This stage is implemented but not currently deployed.
The pipeline terminates at VadcopFinal (Stage~5) in the current configuration.}

\medskip

The final stage would produce a constant-size proof verifiable on Ethereum.

\begin{enumerate}[label=(\roman*)]
  \item \textbf{System.}
        A BN254 SNARK using the fflonk variant of PLONK,
        which requires only one pairing check for verification.

  \item \textbf{publicsHash.}
        The public inputs are hashed via SHA-256 to produce a single
        $\mathrm{publicsHash}$ value that the on-chain verifier checks.

  \item \textbf{On-chain verifier.}
        A Solidity smart contract that:
        \begin{enumerate}[label=(\alph*)]
          \item Recomputes $\mathrm{publicsHash}$ from the provided public inputs.
          \item Verifies the fflonk proof using the BN254 precompiles
                (ecAdd, ecMul, ecPairing).
        \end{enumerate}

  \item \textbf{Output.}
        A constant-size proof $\pi_{\mathrm{fflonk}}$
        (typically $\sim$768 bytes)
        plus the public inputs.
\end{enumerate}

%=============================================================================
\section{Distributed Proving}
\label{sec:distributed}
%=============================================================================

The VADCOP architecture enables distributed proving:
workers independently compute per-AIR proofs
and contribute to challenge aggregation
without sharing witness data.
The three-phase workflow is:

\begin{enumerate}[label=\textbf{Phase \arabic*:}]
  \item \textbf{Contributions.}
        Each worker computes stage-1 commitments and
        lattice contributions for its assigned AIR instances.
        Contributions are sent to the coordinator.

  \item \textbf{Prove.}
        The coordinator aggregates all contributions
        into the global challenge $\chi$
        and distributes it to all workers.
        Workers generate per-AIR STARK proofs (Basic stage)
        and perform internal Recursive1/Recursive2 aggregation
        for their assigned airgroups.

  \item \textbf{Aggregate.}
        A designated aggregator receives all per-airgroup
        Recursive2 proofs, performs final Recursive2 tree reduction
        (if needed), and generates the VadcopFinal proof.
        In the current deployment this is the terminal stage;
        when \texttt{final\_snark} is enabled,
        RecursiveF and fflonk would follow sequentially.
\end{enumerate}

\paragraph{MPI worker aggregation.}
In the MPI deployment model, each worker process handles a subset
of AIR instances.
The coordinator process (rank 0) performs the aggregation
of contributions and final proof assembly.

\paragraph{Coordinator-mediated multi-machine proving.}
For large computations spanning multiple machines,
the coordinator mediates challenge distribution and proof collection
over the network.
Each machine runs an independent set of workers,
and only the aggregated contributions and final proofs
traverse machine boundaries.
